\begingroup
\small
\setlength{\textfloatsep}{5pt plus 1pt minus 1pt}
\setlength{\floatsep}{4pt plus 1pt minus 1pt}
\setlength{\intextsep}{3pt plus 1pt minus 1pt}
\setlength{\abovecaptionskip}{2pt}
\setlength{\belowcaptionskip}{2pt}

This appendix preserves the protocol information and implementation checks that
are most useful for reproducing the main comparison. The remaining exploratory
tables and traces are retained in the repository's full appendix sources.

\subsection{Protocol and baseline fidelity}

The fixed $5\%$ rate is a declared compression operating point rather than a
development-score maximizer. With the student budget fixed, development
Macro-F1 for the $1\%$, $5\%$, and $10\%$ rates was 51.09, 52.95, and 51.47
on the three-stage ArXiv stream, and 79.28, 79.25, and out of memory
on SportsFit. The $10\%$ SportsFit run therefore does not represent a larger
usable memory under the same hardware limit. Historical memory is kept fixed
after each stage. This is the access policy used in the main evaluation.

CaT and GECC retain their distribution-matching and propagated-feature
clustering cores, respectively, and both attach source text for the common
student. GCond-x uses the official simplified graph convolution and gradient
distance objective with the continual proxy adapted to this access protocol.
SFGC uses its trajectory-matching objective with the reduced teacher and
construction budgets. TAGSAM retains its selector and response-matching core
with 100 outer iterations and 10 response minibatch updates per iteration.
the official configuration uses 5,000 outer iterations and 15 full synthetic
data epochs per iteration, so the two settings are adaptations rather than
official reproductions. STAD changes supervision, synthetic topology, and
continual access. These distinctions define the comparison scope and prevent
adapted rows from being read as native reproductions.

The ArXiv comparison is repeated below for auditability because it contains
the larger baseline set used in the main table. It uses the same three-stage
13/14/13 schedule and is not an additional five-stage test. The $1\%$, $5\%$,
and $10\%$ rate check is retained as a development diagnostic. The fixed
$5\%$ rate is the common operating point used in the main tables.

\begin{table}[H]
\centering
\footnotesize
\caption{ArXiv comparison repeated for auditability. The three-stage class schedule is 13/14/13, matching the ArXiv column in \autoref{tab:main_comparison}. Entries are mean $\pm$ SD across nine memories. Each memory is evaluated with three student seeds.}
\label{tab:arxiv_three_stage}
\begin{tabular*}{\linewidth}{@{\extracolsep{\fill}}lcr@{}}
\toprule
Method & $n$ & Macro-F1 / Accuracy\\
\midrule
CaT & 9 & 34.98$\pm$1.79 / 60.41$\pm$2.64\\
GECC & 9 & 52.68$\pm$0.45 / 70.94$\pm$0.49\\
Random & 9 & 47.72$\pm$0.56 / 69.39$\pm$0.30\\
K-center & 9 & 30.84$\pm$1.50 / 59.24$\pm$1.26\\
Cluster medoid & 9 & 48.30$\pm$0.86 / 69.41$\pm$0.27\\
Herding & 9 & 49.21$\pm$0.18 / 69.63$\pm$0.34\\
GCond-x & 9 & 49.64$\pm$1.28 / 66.71$\pm$0.93\\
SFGC & 9 & 24.23$\pm$2.17 / 49.10$\pm$3.41\\
TAGSAM & 9 & 39.31$\pm$3.11 / 58.92$\pm$3.54\\
STAD & 9 & 49.76$\pm$0.64 / 66.07$\pm$1.37\\
\method{} & 9 & 52.91$\pm$0.73 / 71.15$\pm$0.59\\
\bottomrule
\end{tabular*}
\end{table}


\begin{table}[H]
\centering
\scriptsize
\setlength{\tabcolsep}{3pt}
\caption{Archived condensation-rate check on the development partition.
The student budget is fixed. OOM denotes an out-of-memory failure.}
\label{tab:compact_rate}
\begin{tabular*}{\linewidth}{@{\extracolsep{\fill}}lrrr@{}}
\toprule
Dataset & $1\%$ & $5\%$ & $10\%$\\
\midrule
ArXiv (3 stages) & 51.09 & 52.95 & 51.47\\
SportsFit & 79.28 & 79.25 & OOM\\
BookHis & 46.54 & 47.39 & 45.87\\
\bottomrule
\end{tabular*}
\end{table}

\begin{table*}[t]
\centering
\scriptsize
\setlength{\tabcolsep}{4pt}
\caption{Selected protocol and implementation evidence. Development results
are not substituted for the primary test comparison.}
\label{tab:compact_evidence}
\begin{tabular}{@{}p{0.26\textwidth}p{0.69\textwidth}@{}}
\toprule
Evidence & Result or interpretation \\
\midrule
Five-stage ArXiv development check & CaT 35.32$\pm$2.11, GECC 51.31$\pm$0.74, and \method{} 51.93$\pm$1.06 Macro-F1. The paired difference from GECC is +0.62 points with 95\% interval $[-0.45,1.69]$. \\
Optional class-block storage & Peak construction allocation changes from 11.537 to 5.252 GiB on ArXiv and from 23.802 to 17.514 GiB on SportsFit. The corresponding F1 changes are +0.021 and $-0.268$ points. \\
Incremental versus one-call construction & With the optional implementation, peak allocation falls from 8.939 to 5.252 GiB on ArXiv, a 41.25\% reduction, and from 18.851 to 17.514 GiB on SportsFit, a 7.09\% reduction. The comparison is about peak working memory because the total update counts differ. \\
Update sensitivity & Among 65 completed reduced-objective stage traces, 20 satisfy the practical plateau rule of less than 1\% relative absolute loss change in two consecutive 10-update windows. All five BookHis stages satisfy it in the 80--90 and 90--100 windows of the 100-update runs. \\
\bottomrule
\end{tabular}
\end{table*}

\subsection{Chronological data and prediction scope}

The chronological ArXiv stream uses publication-year cutoffs 2015, 2016,
2017, 2018, and 2020. The ACM subset contains 16,928 nodes from databases,
data mining, and networking, with 11,848 construction nodes and 592 final
memory items. The DBLP subset contains 57,520 nodes from artificial
intelligence, computer vision, natural language processing, and machine
learning, with 40,227 construction nodes and 2,011 final memory items. The
training-year quantiles are 2005/2008/2010/2012/2015 for ACM and
2008/2010/2012/2014/2015 for DBLP, and papers from one year remain together.
Five exact normalized-title overlaps with ACM were removed from DBLP. These
are curated chronological subsets rather than unchanged official benchmark
splits.

Each method constructs a new block from the current arrival group and its
induced edges, retains the historical synthetic prefix, and uses the
cumulative $5\%$ budget. The complete 30-update objective is retained. The
chronological runs use float64 Sinkhorn normalization with a 16,384 iteration
cap, tolerance $10^{-4}$, and float32 output. This numerical repair resolves a
failed transport solve without adding outer construction updates. Cumulative
prediction uses all observed test nodes and the graph induced by all arrived
endpoints. A separate current-arrival diagnostic restricts both queries and
inference edges to the current group. It is reported only as a scope check.

\subsection{Resource and convergence checks}

The constructor keeps the full source graph in the CPU loader, while the
logical access rule restricts construction to the current increment and the
historical memory. Dense assignment states and proxy intermediates dominate
the working set, so a sparse edge-free operator alone yields only a small
peak reduction. The optional class-block storage and cache offloading reduce
construction memory without changing the objective, whereas the principal
resource benefit remains downstream reuse of the small memory. Student
training with precomputed text embeddings reduces peak CUDA allocation by
93.8\% on SportsFit and 89.5\% on BookHis relative to training on all
observed construction nodes.

\begin{table}[H]
\centering
\scriptsize
\setlength{\tabcolsep}{2pt}
\caption{Downstream student training. Values average three seeds over 2,048 updates. Text embeddings are precomputed and the encoder is off the GPU.}
\label{tab:student_resource_audit}
\begin{tabular}{cccc}
\toprule
Dataset & 5\% Memory & Full raw graph & Reduction\\
 & MiB / s & MiB / s & Memory / Time\\
\midrule
ArXiv & 100.7 / 7.75 & 1614.6 / 63.61 & 93.8\% / 87.8\%\\
SportsFit & 102.2 / 7.73 & 1656.3 / 66.64 & 93.8\% / 88.4\%\\
BookHis & 43.1 / 7.42 & 411.2 / 15.11 & 89.5\% / 50.9\%\\
\bottomrule
\end{tabular}
\end{table}

The update-sensitivity diagnostic removes GM because the complete objective
requires second derivatives and can exceed available GPU memory. It therefore
describes the reduced objective rather than proving convergence of the full
method. Increasing the reduced objective from 30 to 60 updates changes F1 by
$-0.142$ points on ArXiv and $-2.249$ points on BookHis. Increasing from 30 to
100 changes BookHis F1 by $-1.578$ points. These results support the use of 30
updates as the fixed main-protocol budget while keeping the plateau analysis
separate from downstream accuracy selection.

For completeness, the two resource audits below preserve the earlier
implementation evidence in a compact form. The first compares the dense and
class-block storage variants. The second compares incremental construction
with one call over all observed nodes. They measure peak construction
allocation, not downstream training cost, and their update counts are not
identical.

\begin{table}[H]
\centering
\scriptsize
\setlength{\tabcolsep}{2.5pt}
\caption{Archived construction-memory audits. Entries are peak allocated
GiB. Class blocks include cache offloading.}
\label{tab:compact_memory}
\begin{tabular*}{\linewidth}{@{\extracolsep{\fill}}lrrr@{}}
\toprule
Audit & ArXiv & SportsFit & BookHis\\
\midrule
Dense storage & 11.537 & 23.802 & --\\
Class blocks & 5.252 & 17.514 & --\\
Incremental construction & 5.252 & 17.514 & --\\
One-call full history & 8.939 & 18.851 & --\\
\bottomrule
\end{tabular*}
\end{table}

The corresponding reductions are 54.48\% and 26.42\% for class-block
storage, and 41.25\% and 7.09\% for incremental construction on ArXiv and
SportsFit, respectively. The complete resource tables, text controls,
architecture transfer study, and loss traces remain in the repository's
archived appendix sources for auditability.

\subsection{Additional archived controls}

The following compact tables retain scope and transfer checks from the earlier
appendix. They are not used to select the main protocol. ``Current'' evaluation
restricts both the query nodes and inference edges to the last-stage arrivals.
cumulative evaluation uses the graph observed up to that stage.

\begin{table}[H]
\centering
\scriptsize
\setlength{\tabcolsep}{3pt}
\caption{Current-arrival evaluation from the chronological audit. Entries are
final-stage Macro-F1.}
\label{tab:compact_current}
\begin{tabular*}{\linewidth}{@{\extracolsep{\fill}}lrrr@{}}
\toprule
Method & ArXiv & ACM & DBLP\\
\midrule
CaT & 37.46 & 85.14 & 76.59\\
GECC & 48.10 & 85.10 & 75.30\\
\method{} & 49.83 & 85.18 & 73.75\\
\bottomrule
\end{tabular*}
\end{table}

\begin{table}[H]
\centering
\scriptsize
\setlength{\tabcolsep}{2.5pt}
\caption{Mechanism and assignment controls on the archived three-stage ArXiv
development batch. Cost is peak GiB / seconds.}
\label{tab:compact_controls}
\begin{tabular*}{\linewidth}{@{\extracolsep{\fill}}lrrr@{}}
\toprule
Control & F1 & Accuracy & Cost\\
\midrule
\method{} & 52.96 & 70.64 & --\\
Fixed $Q$ & 40.20 & 55.80 & --\\
Without DM & 51.05 & 64.19 & --\\
Without GM & 53.20 & 70.83 & --\\
Separate assignment & 53.20 & -- & 15.99/292.5\\
Shared assignment & 53.24 & -- & 12.08/241.3\\
\bottomrule
\end{tabular*}
\end{table}

\begin{table}[H]
\centering
\scriptsize
\setlength{\tabcolsep}{2.5pt}
\caption{Archived update sensitivity without GM. Entries are Macro-F1. OOM and
failed denote memory and numerical failures, respectively.}
\label{tab:compact_steps}
\begin{tabular*}{\linewidth}{@{\extracolsep{\fill}}lrrr@{}}
\toprule
Dataset & 30 updates & 60 updates & 100 updates\\
\midrule
ArXiv & 51.702 & 51.809 & failed\\
SportsFit & 79.980 & OOM & OOM\\
BookHis & 48.020 & 45.685 & 46.356\\
\bottomrule
\end{tabular*}
\end{table}

\begin{table}[H]
\centering
\scriptsize
\setlength{\tabcolsep}{2.5pt}
\caption{Additional evaluation partition. Entries are Macro-F1 averaged over
three construction seeds and three student seeds per memory.}
\label{tab:compact_partition}
\begin{tabular*}{\linewidth}{@{\extracolsep{\fill}}lrrr@{}}
\toprule
Dataset & CaT & GECC & \method{}\\
\midrule
ArXiv & 35.70$\pm$0.62 & 51.89$\pm$0.63 & 53.40$\pm$0.13\\
SportsFit & 75.04$\pm$1.90 & 80.30$\pm$0.55 & 79.56$\pm$0.24\\
BookHis & 26.87$\pm$0.99 & 42.01$\pm$0.24 & 45.16$\pm$2.03\\
\bottomrule
\end{tabular*}
\end{table}

\begin{table}[H]
\centering
\scriptsize
\setlength{\tabcolsep}{2.5pt}
\caption{Archived architecture transfer on three-stage ArXiv. Entries are
Macro-F1 / Accuracy on fixed memories.}
\label{tab:compact_architecture}
\begin{tabular*}{\linewidth}{@{\extracolsep{\fill}}lrr@{}}
\toprule
Method & GCN & GraphSAGE\\
\midrule
CaT & 34.98 / 60.41 & 50.24 / 70.78\\
GECC & 52.68 / 70.94 & 52.60 / 70.62\\
\method{} & 52.91 / 71.15 & 53.93 / 71.27\\
\bottomrule
\end{tabular*}
\end{table}

\begin{table}[H]
\centering
\scriptsize
\setlength{\tabcolsep}{2pt}
\caption{Archived support and text controls. Support diagnostics average
recorded stages and seeds. The remaining rows report Macro-F1.}
\label{tab:compact_text}
\begin{tabular*}{\linewidth}{@{\extracolsep{\fill}}lrrr@{}}
\toprule
Diagnostic or control & ArXiv & SportsFit & BookHis\\
\midrule
Hard slot coverage & 1.0000 & 0.9998 & 1.0000\\
Normalized row entropy & 0.1200 & 0.2385 & 0.0827\\
Unique text ratio & 1.0000 & 0.9995 & 0.9990\\
Pairwise text cosine & 0.0947 & 0.3988 & 0.1194\\
Soft/hard cosine gap & 0.3853 & 0.2037 & 0.4693\\
\midrule
Primary text control & 52.96 & 80.02 & 48.53\\
No text initialization & 53.52 & 77.50 & 46.79\\
No pair feedback & 53.13 & 80.30 & 48.14\\
Both off & 53.32 & 77.64 & 47.87\\
\midrule
Nearest-1 text & 53.14 & 79.68 & 47.86\\
Nearest-5 texts & 52.59 & 80.75 & 47.98\\
Random-1 text & 53.23 & 80.33 & 47.56\\
Random-5 texts & 53.01 & 80.65 & 46.56\\
\bottomrule
\end{tabular*}
\end{table}

\endgroup
